% Incorporación al final de extension.tex; no contiene una sección principal.
% Procedencia: RESULTADO_RECUPERADO; explicitación sectorial de las hipótesis.
% Fuentes leídas: monografía Doble proyección holográfica (cantera editorial),
% manuscrito/local/01_tesis_recuperada.tex y
% manuscrito/generated/algebra_holonomica.tex, especialmente §§1--7 y 21.
% Definición rectora de número: manuscrito/flattened/main_autosuficiente.tex,
% secciones que comienzan en líneas 22137 y 23000 de la cantera de 775 páginas.
% No se adopta su presentación global como prueba de todas las flechas TPK.

\subsection{Le nombre comme section compatible et la valeur comme lecture}
\label{hist:seccion}

La distinction entre état et observation permet de préciser quel objet se
prolonge lorsque la résolution augmente. À la profondeur \(n\), soit \(X_n\)
l’ensemble des états admissibles produits par APP--TRIT--TPK. Leurs données
comprennent la position, les feuillets somme et produit, les résidus et quotients, le régime,
l’orientation, la retenue, la phase, la route, la frontière, la mémoire et la survie. Ce ne sont pas
des coordonnées indépendantes : les évaluations arithmétiques et les transports
précédents imposent leurs relations. Un prolongement conserve ces relations
dans chaque antécédent, même s’il ajoute de nouvelles étapes et de nouvelles données de mémoire.
La définition suivante formalise cette arithmétique des histoires au moyen
de ses applications de restriction.

Soient \(r_{m,n}:X_m\to X_n\), pour \(m\geq n\), les restrictions de
profondeur, avec \(r_{n,n}=\operatorname{id}\) et
\(r_{\ell,n}=r_{m,n}\circ r_{\ell,m}\). Chaque restriction conserve le
préfixe et l’information déjà déterminée en lui.

\begin{definition}[Section numérique compatible]
Un nombre HMT est une histoire compatible du système d’états :
\begin{equation}
 x=(x_n)_{n\geq0}\in X_\infty:=\varprojlim_n X_n,
 \qquad r_{m,n}(x_m)=x_n.
 \label{hist:limite-inverso}
\end{equation}
Un lecteur est une opération supplémentaire sur un domaine spécifié de ces
histoires. Un chiffre est une observation finie ; une valeur scalaire est une
évaluation du lecteur. Le couple formé de la section et de son lecteur constitue
une présentation munie d’une lecture, non une redéfinition de l’identité de la section.
\end{definition}

Dans le domaine archimédien \(X_\infty^{\mathrm{Arch}}\), le lecteur associe
à chaque préfixe un cylindre fermé, conserve leur emboîtement et fait tendre
leur diamètre vers zéro. Il détermine ainsi
\begin{equation}
 \operatorname{ev}_{\mathbb R}:X_\infty^{\mathrm{Arch}}\longrightarrow
 \mathbb R,\qquad
 \{\operatorname{ev}_{\mathbb R}(x)\}=\bigcap_n I_n(x).
 \label{hist:evaluacion}
\end{equation}
La construction concrète de ces cylindres et la décision de leurs frontières
sont développées dans la section~\ref{sec:generacion}. La définition
\eqref{hist:limite-inverso} n’affirme pas à elle seule que l’image de
\eqref{hist:evaluacion} soit la droite entière : une affirmation de couverture
requiert son propre théorème. Une égalité de valeurs n’établit pas non plus, sans
résultat d’injectivité, l’égalité de leurs histoires. Cette séparation
permet de parler de précision comme profondeur en conservant la différence
entre l’antécédent arithmétique et sa réalisation scalaire.

\subsection{Composition des histoires et multiplicateur de mémoire}
\label{hist:algebra}

Fixons un secteur d’états à la profondeur \(n\) et une catégorie
\(\mathcal C\) dont les flèches sont leurs histoires admissibles. La composition
s’écrit \(vu=v\circ u\) : on parcourt d’abord \(u\), puis \(v\).
Les pas élémentaires proviennent de la sélection, du transport et de
l’actualisation du TPK. Une fibre tritique locale admet l’écriture
\(\mathbb R[J_x]/(J_x^2+\tau(x)1_x)\), où
\(\tau(x)\in\{+1,0,-1\}\). Le régime et l’orientation appartiennent à
l’état transporté ; une transition entre régimes n’équivaut pas à
l’identification de leurs algèbres quadratiques.

La formulation suivante précise un secteur algébrique de cette composition.
Pour chaque objet \(x\), on se donne une algèbre réelle associative unitaire \(A_x\),
éventuellement non commutative. Pour chaque \(u:x\to y\), on se donne un
homomorphisme unital \(\rho_u:A_x\to A_y\). Nous supposons une action stricte :
\begin{equation}
 \rho_{\operatorname{id}_x}=\operatorname{id}_{A_x},\qquad
 \rho_v\rho_u=\rho_{vu}.
 \label{hist:accion-estricta}
\end{equation}
Pour chaque couple composable \(u:x\to y\), \(v:y\to z\), nous fixons un
multiplicateur central inversible
\(\Omega(v,u)\in Z(A_z)^\times\). On exige la normalisation
\(\Omega(\operatorname{id}_y,u)=\Omega(u,\operatorname{id}_x)=1_y\)
et, pour chaque triplet composable, l’identité
\begin{equation}
 \rho_w\bigl(\Omega(v,u)\bigr)\Omega(w,vu)
   =\Omega(w,v)\Omega(wv,u).
 \label{hist:cociclo}
\end{equation}
Il s’agit d’hypothèses explicites sur les transports et leurs coefficients, non d’une
conséquence du qualificatif holonomique appliqué à la dynamique. Pour appliquer cette présentation
à un secteur TPK, il faut identifier ses \(\rho_u\) et \(\Omega(v,u)\).
Les transitions exprimées par des bimodules exigent la conservation de cette
structure et ne sont pas remplacées par \eqref{hist:accion-estricta}.

L’espace des sommes finies
\[
 \mathcal H(\mathcal C,A,\rho,\Omega)
   =\bigoplus_{u\in\operatorname{Mor}(\mathcal C)} A_{t(u)}T_u
\]
est muni du produit bilinéaire
\begin{equation}
 (aT_v)(bT_u)=a\rho_v(b)\Omega(v,u)T_{vu},
 \label{hist:producto}
\end{equation}
si les flèches sont composables, et de zéro sinon. Ici, \(t(u)\)
désigne l’état terminal. La somme réunit des histoires à coefficients ; le
produit concatène les histoires en conservant le multiplicateur du registre.

\begin{proposition}[Associativité sectorielle à multiplicateur central]
Sous \eqref{hist:accion-estricta} et \eqref{hist:cociclo}, le produit
\eqref{hist:producto} est associatif. Si l’ensemble des objets est fini,
son unité est \(\sum_x1_xT_{\operatorname{id}_x}\).
\end{proposition}
\begin{proof}
Pour \(u:x\to y\), \(v:y\to z\), \(w:z\to t\), soient
\(a\in A_y\), \(b\in A_z\), \(c\in A_t\). Les deux coefficients
de \(T_{wvu}\) sont
\begin{align*}
 ((cT_w)(bT_v))(aT_u):\quad&
 c\rho_w(b)\Omega(w,v)\rho_{wv}(a)\Omega(wv,u),\\
 (cT_w)((bT_v)(aT_u)):\quad&
 c\rho_w(b)\rho_w\rho_v(a)\rho_w(\Omega(v,u))\Omega(w,vu).
\end{align*}
L’action stricte identifie \(\rho_w\rho_v(a)\) à \(\rho_{wv}(a)\).
La centralité de \(\Omega(w,v)\) permet de le déplacer à droite de
ce facteur ; \eqref{hist:cociclo} égalise alors les produits restants.
Les coefficients \(a,b,c\) ne sont pas permutés. Si le triplet n’est pas composable,
les deux parenthésages sont nuls du fait de leurs états initiaux et terminaux.
La bilinéarité achève la preuve. Les identités locales et la
normalisation de \(\Omega\) vérifient l’unité indiquée.
\end{proof}

La retenue du calendrier fournit un exemple effectif. Dans la catégorie
à un objet et de flèches \(C_9\), prenons pour coefficients
\(A=\mathbb R[z,z^{-1}]\), l’action identité et
\[
 \Omega(b,a)=z^{c_0(a,b)},\qquad
 c_0(a,b)=\left\lfloor\frac{a+b}{9}\right\rfloor,
 \quad 0\leq a,b<9.
\]
Les deux sommes de retenues d’un triplet sont
\(\lfloor(a+b+c)/9\rfloor\) ; \eqref{hist:cociclo}
est donc satisfaite. La variable formelle \(z\) enregistre le tour sans
lui attribuer de valeur numérique. Imposer ensuite \(z^{12}=1\) restitue le
registre dodécaphasique du calendrier \eqref{eq:extension108} ; conserver
\(z\) sans cette relation distingue les tours successifs. L’exemple réalise
la mémoire de ce calendrier, sans l’identifier à la totalité de la
mémoire TPK.

\subsection{Double variance du raffinement et conservation de l’unité}
\label{hist:doble-varianza}

La restriction des états prolonge les observables en sens contraire.
Considérons les algèbres de fonctions, munies du produit ponctuel,
\(D_n=\operatorname{Fun}(X_n,\mathbb R)\). La précomposition définit
\begin{equation}
 r_{m,n}^*:D_n\longrightarrow D_m,
 \qquad r_{m,n}^*f=f\circ r_{m,n}.
 \label{hist:precomposicion}
\end{equation}
Ce sont des morphismes unitaux ; ils sont également injectifs lorsque les restrictions
correspondantes sont surjectives. La limite inductive
\(D_{\mathrm{cyl}}=\varinjlim_n(D_n,r_{n+1,n}^*)\) réunit les
observations de profondeur finie. Cette algèbre n’est pas identifiée ici
à l’algèbre entière des routes : prolonger les opérateurs de transport exige
également leurs compatibilités avec le raffinement.

\begin{lemma}[Unité et naturalité de l’évaluation]
Si \(e_x\) est l’indicatrice de \(x\in X_n\), on a
\begin{equation}
 r_{m,n}^*e_x=\sum_{y:\,r_{m,n}(y)=x}e_y,
 \qquad r_{m,n}^*1_n=1_m,
 \label{hist:unidad}
\end{equation}
et, pour tous \(f\in D_n\), \(y\in X_m\),
\[
 \langle r_{m,n}^*f,y\rangle=\langle f,r_{m,n}y\rangle.
\]
Chaque section \(x\in X_\infty\) détermine une évaluation bien définie
\(D_{\mathrm{cyl}}\to\mathbb R\), donnée par \([f]\mapsto f(x_n)\)
pour \(f\in D_n\).
\end{lemma}
\begin{proof}
Les fibres de \(r_{m,n}\) forment une partition de \(X_m\). Leurs
indicatrices sont orthogonales et de somme \(1_m\), ce qui démontre
\eqref{hist:unidad}. La naturalité est la définition de la précomposition.
Enfin, \((r_{m,n}^*f)(x_m)=f(x_n)\), de sorte que deux représentants
compatibles donnent la même évaluation dans la limite inductive.
\end{proof}

L’unité globale est la fonction constante un sur toutes les possibilités ;
la section individuelle sélectionne une histoire compatible. Conserver la
première ne signifie pas immobiliser la seconde. Le raffinement peut augmenter
le nombre d’extensions et la profondeur de chaque histoire en maintenant
l’unité et toutes les évaluations de ses antécédents.

\subsection{Retour de phase, représentation de Weyl et deux lectures}
\label{hist:puente-lecturas}

La loi \eqref{eq:holonomia-memoria} acquiert une réalisation opératorielle
dans la suspension bilatérale qui sera construite dans la
section~\ref{sec:generacion}. Dans celle-ci, \(T\) est l’avancement d’un pas,
\(V_9\) observe la phase nonadique, \(W_\delta\) la phase de rotation et
\(D_M\) le degré de mémoire. En écrivant
\(\widehat\Gamma_9=T^9\) pour le tour dans cette représentation,
les relations de Weyl \eqref{eq:weyl-relojes} impliquent, sur les vecteurs
de support fini,
\begin{equation}
 \begin{aligned}
 V_9\widehat\Gamma_9&=\widehat\Gamma_9V_9,\qquad
 W_\delta\widehat\Gamma_9=e^{2\pi i9\delta}
     \widehat\Gamma_9W_\delta,\\
 [D_M,\widehat\Gamma_9]&=\widehat\Gamma_9.
 \end{aligned}
 \label{hist:tres-retornos}
\end{equation}
La première égalité découle de \(\zeta_9^9=1\) ; la deuxième, de neuf
applications de la covariance de \(W_\delta\) ; la troisième exprime
l’incrément unitaire de mémoire. La phase finie revient, tandis que la
rotation irrationnelle et la mémoire distinguent l’état suivant. Les
caractères et le paramètre \(\delta\) sont réalisés ultérieurement à partir de la
construction numérique spécifiée dans cette section : ils ne sélectionnent pas rétrospectivement
les conditions initiales. Cette représentation décrit les observables et les translations de
la suspension ; elle ne remplace ni toutes les coordonnées ni tous les opérateurs
du TPK.

La double lecture prolonge la même distinction entre histoire et observable.
Dans son domaine archimédien, on applique \eqref{hist:evaluacion} ; dans le
domaine d’incidence calibré, les applications \(Q_n\) de la
section~\ref{exc:doble-lectura} conservent les mots et les repères de
frontière. L’égalité qui y est démontrée,
\(r_{n,m}Q_n=Q_mp_{n,m}\), détermine leur prolongement \(Q_\infty\).
Sur le domaine commun aux deux lectures, on peut considérer conjointement
\(\operatorname{ev}_{\mathbb R}(x)\) et \(Q_\infty(x)\) : la
première retient la valeur ; la seconde, l’incidence transportée que son
codomaine conserve. La section précède les deux. Cette articulation permet
de passer de l’arithmétique des histoires aux constantes et à la lecture
exceptionnelle sans transformer leurs réalisations en nouvelles données initiales.
